% T6_Potentials.pdf -- Cvičení: termodynamické potenciály
\section{Cvičení: Termodynamické potenciály}
\strana{T6-1}

\begin{priklad}[Příklad 1]
Najděte Legendreovu transformaci fce $y(x) = Ae^{Bx}$
\hfill \textcolor{NavyBlue}{$\psi(p) = \frac{p}{B}\left(1 - \log\frac{p}{AB}\right)$}
\end{priklad}

\begin{priklad}[Příklad 2]
$F, G, U, H$ pro jednokomponentní ideální plyn (vše jako FR v přirozených proměnných)
\end{priklad}

\color{NavyBlue}
\underline{Řešení:} \qquad $pV = Nk_BT$ \quad \& \quad $U = cNk_BT$
\begin{enumerate}[label=\arabic*)]
  \item $\vzorec{F = Nk_BT\left\{c - \dfrac{s_0}{k_B} - \log\left[\left(\dfrac{T}{T_0}\right)^c\left(\dfrac{V}{V_0}\right)\left(\dfrac{N_0}{N}\right)\right]\right\}}$
  \begin{itemize}
    \item bylo na přednášce minulý týden (vyšli jsme ze znalosti entropické FR a kalorické stav.\ rovnice)
    \item alternativně lze nalézt přímo řešením parciální rovnice
    \[ F(T,V,N) = U(T) + T\left(\pdv{F}{T}\right)_{V,N} \]
    ovšem integrační konstanta nebude vztažena k $s_0$
  \end{itemize}
  \item $\vzorec{G(T,p,N) = Nk_BT\left(c + \dfrac{s_0}{k_B} - \log\left[\left(\dfrac{T}{T_0}\right)^{c+1}\left(\dfrac{p_0}{p}\right)\right] + 1\right)}$
  \begin{itemize}
    \item triviální LT \quad $G = F(T,p,N) + pV(T,p,N)$, potřebujeme pouze termickou stav.\ rovnici
  \end{itemize}
  \pozn{znaménko u $s_0/k_B$ se v originále liší mezi $F$ a $G$; přepsáno jak zapsáno}
  \item $U(S,V,N) = $ \uv{inverze entropické FR}
  \item $\vzorec{H = \dfrac{c+1}{c}U(S,p,N) = \dfrac{c+1}{c}U_0\left(\dfrac{N}{N_0}\right)\left(\dfrac{p}{p_0}\right)^{\frac{1}{c+1}}\exp\left(\dfrac{S - Ns_0}{(c+1)Nk_B}\right)}$ \\
  \hspace*{0.6cm}$\uparrow$ \ $U = cpV$ \quad \& \quad $H = U + pV$
\end{enumerate}
\color{black}

%----------------------------------------------------------------------
\strana{T6-2}
\begin{priklad}[Příklad 3]
$U, F, G, H$ pro fotonový plyn \qquad $U = bVT^4$ \quad \& \quad $p = \dfrac{U}{3V}$
\end{priklad}

\color{NavyBlue}
\underline{Řešení:} (bez záruky)
\begin{enumerate}[label=\arabic*)]
  \item $S = \frac43 b^{1/4}U^{3/4}V^{1/4}$ \quad -- bylo někdy na přednášce, $N$ není TD proměnná (mikroskop.\ se nezachovává)
  \item $U(S,V) = \left(\frac34\right)^{4/3}b^{-1/3}S^{4/3}V^{-1/3}$ \quad -- jen inverze
  \item $\vzorec{F(T,V) = -\frac13 bVT^4}$ \qquad $[= U(T,V) - TS(T,V)]$
  \item $G(T,p) = F(T,p) + pV(T,p) = 0$ \qquad (nezbyla ext.\ proměnná)
  \item $H(S,p) = \left(\frac3b\right)^{1/4}Sp^{1/4}$ \qquad $[= 4pV(S,p)]$
\end{enumerate}
\color{black}

\begin{priklad}[Příklad 4*]
TD potenciály pro van der Waalsův plyn
\end{priklad}
\textcolor{NavyBlue}{-- spíše pro otrlé, $G(T,p,N)$ myslím ani nelze explicitně vyjádřit kvůli inverzi $V = V(T,p,N)$}

\begin{priklad}[Příklad 5*]
Ukažte, že pro multikomponentní id.\ plyn je volná energie aditivní přes komponenty,
\[ F(T,V,N_1,\dots,N_\ell) = \sum_{i=1}^\ell F(T,V,N_i), \]
a že podobná aditivita neplatí pro ostatní potenciály.
\end{priklad}

